\chapter*{Lời Mở Đầu}
\addcontentsline{toc}{chapter}{Lời Mở Đầu}
\markboth{Lời Mở Đầu}{Lời Mở Đầu}

\begin{truyen}{Lời Mở Đầu}{Opening Words}
\chuhoa{H}{ọc trò có rất nhiều câu hỏi mà không dám hỏi thẳng.}
\textit{(Students have many questions they dare not ask directly.)}

Hỏi thẳng có nghĩa là phải thừa nhận mình chưa biết. Phải để lộ điều mình đang lo lắng. Phải chấp nhận rằng câu trả lời có thể không như mình muốn nghe. Đó là những lý do khiến nhiều câu hỏi quan trọng nhất bị giữ lại trong im lặng.
\textit{(Asking directly means admitting you do not know. It means revealing what worries you. It means accepting that the answer may not be what you want to hear. These are the reasons the most important questions are kept in silence.)}

Quyển mười ba này chọn những câu hỏi đó ra — và trả lời chúng một cách thẳng thắn, không vòng vo.
\textit{(Volume thirteen selects those questions --- and answers them directly, without detour.)}

Không phải mọi câu trả lời đều dễ nghe. Nhưng tất cả đều thật.
\textit{(Not every answer is easy to hear. But all of them are honest.)}
\end{truyen}

\begin{baihoc}
Câu hỏi không được hỏi là câu hỏi không bao giờ được trả lời --- và sự im lặng đó tốn nhiều hơn ta nghĩ.
\textit{(The question never asked is the question never answered --- and that silence costs more than we think.)}
\end{baihoc}

\begin{ghinhoanh}
\textbf{Short English to remember:}

Unasked questions stay in the dark.

Honest answers are the most respectful kind.
\end{ghinhoanh}

\begin{tuhocnhanh}
1. Câu hỏi nào em đang giữ trong im lặng mà chưa dám hỏi ai? \textit{What question are you currently keeping in silence, afraid to ask?}

2. Điều gì ngăn em đặt câu hỏi thẳng thắn với thầy/cô? \textit{What stops you from asking your teacher direct questions?}

3. Nếu em có thể hỏi thầy/cô một câu và được trả lời thật thật thà, em hỏi gì? \textit{If you could ask your teacher one question and receive a completely honest answer, what would you ask?}
\end{tuhocnhanh}

\ngancach
